\begin{abstract}
\textbf{Accuracy on an instructed-lie benchmark does not identify deception-specific signal.} A detector can score high on one while measuring whether the model was \emph{instructed} to lie, and the cause is structural: the instruction to deceive is also the intervention, producing both the deception and the compliance, so the two causal routes are observationally equivalent under \emph{any} manipulation of it. Such a benchmark identifies $\tau_E$, the effect of the \emph{instruction} on the detector, whereas every deception-detection claim is about $\tau_D$, the effect of \emph{deception}---and $\tau_E$ neither bounds nor signs $\tau_D$. \textbf{No increase in accuracy, model scale, or detector sophistication can resolve this unless the evaluation introduces an intervention that separates deception from instruction-following.} Reimplemented from scratch, the standard black-box behavioral detector reaches 81.7\% on instructed-lie trials and \textbf{53.7\%---chance---}once both conditions receive identical prompts; our own five-feature pipeline falls from 93.9--100\% to 61--84\% under that same control. We therefore contribute a \textbf{causal design standard} for such benchmarks, operationalized as an \textbf{audit protocol}: five criteria required before accuracy licenses a claim about deception, four testable on any published benchmark without new data and exactly one able to identify $\tau_D$---and a widely used instructed-roleplay paradigm fails it.

A pre-registered replication on a disjoint claim set and an independently worded probe bank \textbf{eliminates above-chance discrimination for all six targets} (97.0\%$\to$\textbf{43.3\%}), so the collapse holds on \emph{both} claim sets---two mechanistically unrelated detectors failing under one control. \emph{Within} one instructed-deception condition---one deception label, held fixed---the detector still moves $-1.37$ and $+1.61$\,SD in \emph{opposite} directions ($p\!=\!0.0001$), as large as the instruction's own effect, so the score is not a function of the label a benchmark grades against. A parameter-free lexical rule reaches \textbf{69--80\% depending on the pattern inventory}---a lower bound on how much accuracy needs no deception-specific feature. Finally we audit \textbf{ten published rollout sets} for the one criterion that can attribute a signal to deception, namely deception varying at \emph{fixed} elicitation: \textbf{nine cannot supply it}, and on the one that can, the rule is \textbf{null} ($+4.1$/$+7.7$\,pp, $p\!\geq\!0.17$). Scoring that benchmark's hidden scaffold instead would have reported $+84.3$\,pp---the model echoing ``incorrect'' from its own instruction. \textbf{A benchmark must therefore expose only the channel a deployed detector could read}, or hidden prompt leakage is scored as deception signal.

We claim identification, not attribution: these detectors \emph{can} classify prompted deception and we reproduce it, but instructed-roleplay accuracy cannot establish deception as the \emph{cause} of the signal. \textbf{Scope: English, instructed-roleplay evaluations, open-weight models 3B--70B.}
\end{abstract}
